\section{Matching residue characteristics}\label{sec:matching}

We now show that almost every prime of $\Q(i)$ above a split rational prime has an active partner of the same residue characteristic. The auxiliary prime defining activity may vary with the pair. Two finiteness statements make this possible: almost every such target prime has only finitely many possible partners, and the number of target primes missed at a fixed auxiliary prime admits a uniform bound.

Keep the surjection $\beta:G_k\to P$ and the set $\mathcal L$ from Theorem~\ref{thm:normalized-map}, and put
\[
 B=\Q(i),\qquad G_K=\beta^{-1}(P_B),\qquad
 \beta_B=\beta|_{G_K}:G_K\twoheadrightarrow P_B.
\]
Thus $K/k$ is quadratic. An \emph{active pair at $\ell$} in this section means a source prime $v$ of $K$ active at $\ell$ for $\beta_B$, together with its corresponding prime $w$ of $B$ from Theorem~\ref{thm:local-correspondence}. In particular, both primes lie away from $\ell$. For a prime $w$ of $B$, define
\[
 \mathcal V_w=
 \{v:\text{$(v,w)$ is an active pair at some $\ell\in\mathcal L$}\}.
\]

\subsection{Finite fibers and a uniform bound on missed primes}

\begin{lemma}\label{lem:finite-fibers}
For all but finitely many primes $w$ of $B$ above rational primes split in $B$, the set $\mathcal V_w$ is finite.
\end{lemma}
\begin{proof}
Fix $r\in\mathcal L$. Proposition~\ref{prop:signed-representation} supplies a nonzero character $\psi\in\X{r}{B}$ on which $\Gal(B/\Q)$ acts by its nontrivial character. By Proposition~\ref{prop:character-detection}(i), $\psi(\Fr_w)$ is nonzero at all but finitely many degree-one primes $w\nmid r$. Fix such a $w$ above a split rational prime $q\ne r$, and put
\[
 x_0=(\log\chi_r)|_{P_B},\qquad
 z=\frac{\psi(\Fr_w)}{\log_r q}.
\]
The denominator is nonzero: the kernel of the logarithm on $\Z_r^\times$ consists of roots of unity, and the positive rational integer $q$ is not a root of unity. Thus $z\ne0$, and $\psi-zx_0$ vanishes on the decomposition group at $w$, since it is unramified there and vanishes on Frobenius.

Consider its pullback
\[
 y=\beta_B^*(\psi-zx_0)\in\X{r}{K}.
\]
The pullback of $\psi$ transforms by sign under $G_k/G_K$. Its corestriction to $k$ is therefore zero: restricting that corestriction back to $K$ gives the sum of the two translates, and restriction is injective for additive characters with characteristic-zero coefficients. Theorem~\ref{thm:normalized-map} identifies $\beta_B^*x_0$ with $\log\chi_{K,r}$. Consequently
\[
 \cor_{K/\Q}(y)
 =-z[K:\Q]\log\chi_{\Q,r}\ne0.
\]
For every $v\in\mathcal V_w$ away from $r$, decomposition containment gives $y(\Fr_v)=0$, regardless of which auxiliary prime makes the pair active. Proposition~\ref{prop:character-detection}(iii) permits only finitely many such $v$. There are also only finitely many primes of $K$ above $r$, proving the assertion.
\end{proof}

For $\ell\in\mathcal L$, let $\mathcal M_\ell$ be the set of primes of $B$ which occur in no active pair at $\ell$. This includes the primes above $\ell$.

We seek a bound for $\#\mathcal M_\ell$ depending only on $[K:\Q]$. The proof will place independent classes attached to missed primes in one character space of an $S$-unit group. We first record the standard cohomological identification of that group, including the completion step; see also \cite{neukirch1999}.

\begin{lemma}[Unramified Kummer classes]\label{lem:kummer-unramified}
Let $F$ be a number field, let $\ell$ be a rational prime, and let $S$ be a finite set of places of $F$ containing all places above $\ell$ and all archimedean places. Then Kummer theory induces a natural isomorphism
\begin{equation}\label{eq:unramified-kummer}
 H^1_{\mathrm{cont}}(G_F,\Q_\ell(1))
 _{\text{unramified outside }S}
 \simeq \mathcal O_{F,S}^{\times}\otimes_{\Z}\Q_\ell.
\end{equation}
It is equivariant for every group of field automorphisms of $F$ preserving $S$, with the natural action on cohomology combining conjugation on $G_F$ and the action on the Tate twist.
\end{lemma}
\begin{proof}
Integral Kummer theory gives
\[
 H^1_{\mathrm{cont}}(G_F,\Z_\ell(1))
 \simeq\varprojlim_n F^\times/(F^\times)^{\ell^n}.
\]
The inverse-limit assertion follows from the finite-level Kummer isomorphisms; the degree-zero groups are finite and satisfy the Mittag--Leffler condition. Every continuous rational cocycle has bounded image, since $G_F$ is compact, and hence becomes integral after multiplication by one power of $\ell$. The same observation for coboundaries gives rationalization of this integral cohomology. At a prime outside $S$, inertia acts trivially on $\Z_\ell(1)$, and its Kummer evaluation is the valuation times the tame $\ell$-adic character. These integral inertia homomorphisms are torsion-free. Thus unramifiedness after rationalization is exactly the vanishing of all integral valuation coordinates outside $S$.

For completeness, put $U_S=\mathcal O_{F,S}^\times$ and let
\[
 H_n=\{[a]\in F^\times/(F^\times)^{\ell^n}:
       \operatorname{ord}_v(a)\equiv0\pmod{\ell^n}
       \text{ for every }v\notin S\}.
\]
Writing $\Cl_S(F)$ for the $S$-ideal class group, there is an exact sequence
\begin{equation}\label{eq:s-unit-completion}
 0\longrightarrow U_S/U_S^{\ell^n}
 \longrightarrow H_n
 \longrightarrow \Cl_S(F)[\ell^n]
 \longrightarrow0.
\end{equation}
The last map sends $[a]$ to the ideal class of
$\operatorname{div}_S(a)/\ell^n$, where $\operatorname{div}_S$ records the valuations outside $S$. The transition from level $n+1$ to level $n$ induces multiplication by $\ell$ on these obstruction classes. Since $\Cl_S(F)$ is finite, their inverse limit is zero. It follows that the kernel of all outside-$S$ valuations in the completed multiplicative group is precisely $\varprojlim_n U_S/U_S^{\ell^n}$. The argument takes place in the full completed multiplicative group; it does not assume finite valuation support for its elements. Finally $U_S$ is finitely generated, so rationalizing its completion gives $U_S\otimes\Q_\ell$. This proves \eqref{eq:unramified-kummer}. All the constructions commute with field automorphisms preserving $S$, giving the asserted equivariance.
\end{proof}

Uchida's proof of target-side finiteness associates Kummer classes to missed primes and bounds them by a unit rank \cite[Proposition~1]{uchida1981}. We refine that argument by keeping the classes in one character space of a cyclic extension. The multiplicity of this character has a bound independent of $\ell$, even when the extension degree grows.

\begin{proposition}\label{prop:uniform-missed-primes}
For every $\ell\in\mathcal L$,
\[
 \#\{w\in\mathcal M_\ell:w\nmid\ell\}\le 2[K:\Q],
 \qquad
 \#\mathcal M_\ell\le 2[K:\Q]+2.
\]
\end{proposition}
\begin{proof}
Fix $\ell\in\mathcal L$ and any finite set $J\subset\mathcal M_\ell$ of primes away from $\ell$. For each $w\in J$, choose $a_w\in B^\times$ and $h_w>0$ with
\[
 (a_w)=w^{h_w}.
\]
Its compatible Kummer cocycles define a class
\[
 \kappa_w\in H^1_{\mathrm{cont}}(P_B,\Q_\ell(1)).
\]
They factor through $P_B$ because $\Q^s$ contains the roots of unity and all the required radicals. Valuations at the distinct primes $w$ show that these classes are linearly independent. Each is unramified at every finite prime other than $w$ and the primes above $\ell$.

Write
\[
 \rho=\chi_\ell\circ\beta_B,\qquad
 \epsilon=\rho\chi_{K,\ell}^{-1}.
\]
By Theorem~\ref{thm:normalized-map}, $\epsilon$ takes values in $\mu_{\ell-1}\subset\Q_\ell^\times$. Let $K'/K$ be the cyclic extension fixed by its kernel, and let $\Delta=\Gal(K'/K)$. Pullback along the surjection $\beta_B$ is injective on $H^1$, with pulled-back coefficient module $\Q_\ell(\rho)$. Restriction to $G_{K'}$ remains injective: its kernel is a first cohomology group of the finite group $\Delta$, which vanishes in characteristic zero. We therefore obtain independent classes
\[
 c_w\in H^1_{\mathrm{cont}}(G_{K'},\Q_\ell(1))
 \qquad(w\in J).
\]
Under the natural $\Delta$-action on this cohomology, they all lie in the $\epsilon^{-1}$ eigenspace. Indeed the original coefficient module is $\Q_\ell(1)\otimes\epsilon$, so invariance for its twisted action becomes the eigenvalue $\epsilon^{-1}$ after restriction.

We claim that every $c_w$ is unramified away from $\ell$. Suppose that it is ramified at a prime $u\nmid\ell$ of $K'$, and let $v$ be the prime below it in $K$. The cyclotomic module is trivial on $I_u$, so the restricted cocycle is a nonzero continuous homomorphism from $I_u$ to the additive group $\Q_\ell(1)$. Choose a pro-$\ell$ Sylow subgroup $J_u\subset I_u$ on which it is nonzero, and a pro-$\ell$ Sylow subgroup $J_v\subset I_v$ containing $J_u$. Then $\beta_B(J_v)\ne1$, so $v$ is active at $\ell$. Theorem~\ref{thm:local-correspondence} sends $J_v$ into inertia at a corresponding prime $t\nmid\ell$ of $B$. Since $t$ occurs in an active pair, $t\notin J$, whereas $\kappa_w$ is unramified at $t$. Inertia at $t$ acts trivially on the coefficient module, so the Kummer cocycle itself vanishes on that inertia group. Its pullback therefore vanishes on $J_u$, a contradiction. This proves the claim.

Let $S$ consist of the archimedean places of $K'$ and its primes above $\ell$. By Lemma~\ref{lem:kummer-unramified}, the classes $c_w$ give independent elements of the $\epsilon^{-1}$ eigenspace in $\mathcal O_{K',S}^{\times}\otimes\Q_\ell$. It remains to bound this eigenspace. The equivariant Dirichlet $S$-unit theorem \cite[Section~3]{bass1966} gives the character of the full $S$-unit module as
\[
 [\mathcal O_{K',S}^{\times}\otimes\Q_\ell]
   =[\Q_\ell[S]]-[\Q_\ell].
\]
Each $\Delta$-orbit in $S$ contributes at most one copy of the one-dimensional character $\epsilon^{-1}$ to the permutation representation: for a stabilizer $H\subset\Delta$, its multiplicity in $\Q_\ell[\Delta/H]$ is one if it is trivial on $H$, and zero otherwise. There are
\[
 \#(S/\Delta)
 =r_1(K)+r_2(K)+\#\{v\text{ of }K:v\mid\ell\}
 \le2[K:\Q]
\]
such orbits. Subtracting the trivial representation cannot increase this multiplicity. Hence the independent classes $c_w$ lie in a space of dimension at most $2[K:\Q]$, and $\#J\le2[K:\Q]$. Since $J$ was arbitrary, this bounds all missed primes away from $\ell$. There are at most two primes of $B$ above $\ell$, giving the second bound.
\end{proof}

The bound depends on $K$, but neither on $\ell$ nor on the degree of the cyclic extension $K'/K$. We can therefore choose one auxiliary prime that excludes every possible partner of several specified target primes, and contradict this bound. The following elementary Kummer construction provides that choice.

\subsection{Simultaneous separation by power residues}

\begin{lemma}\label{lem:power-separation}
Let $m,d\ge1$, let $q_1,\ldots,q_n$ be distinct rational primes, and for each $j$ let $\mathcal P_j$ be a finite set of rational primes not containing $q_j$. Given any finite set of rational primes to avoid, there are infinitely many odd primes $\ell\equiv1\pmod m$ outside that set and all the displayed prime sets such that
\[
 p^f\bmod\ell\notin\langle q_j\bmod\ell\rangle
 \quad
 (1\le j\le n,\ p\in\mathcal P_j,\ 1\le f\le d).
\]
The sets $\mathcal P_j$ may overlap, and may contain $q_i$ for $i\ne j$.
\end{lemma}
\begin{proof}
Choose distinct auxiliary primes $s_j$ larger than $m$, $d$, and every prime in the lists, and put
\[
 C=\Q(\zeta_{m\prod_j s_j}).
\]
For each $j$, the classes of the distinct rational primes in $\{q_j\}\cup\mathcal P_j$ are independent in $C^\times/(C^\times)^{s_j}$. Indeed, at any prime in those lists, the ramification index in $C/\Q$ divides $[\Q(\zeta_m):\Q]$: the other cyclotomic conductors are prime to that prime. This index is prime to $s_j$, and valuations at the separate rational primes give the asserted independence.

Kummer theory therefore identifies the Galois group over $C$ of
\[
 E_j=C\bigl(q_j^{1/s_j},\ p^{1/s_j}\ (p\in\mathcal P_j)\bigr)
\]
with an elementary abelian $s_j$-group having one independent coordinate for each listed prime. Choose $\sigma_j\in\Gal(E_j/C)$ which fixes $q_j^{1/s_j}$ and acts nontrivially on every $p^{1/s_j}$ with $p\in\mathcal P_j$. The extensions $E_j/C$ are linearly disjoint, since their degrees are powers of distinct primes. Thus these automorphisms combine to an element $\sigma\in\Gal(E/C)$, where $E$ is their compositum.

The extension $E/\Q$ is Galois. Moreover, conjugation in $\Gal(E/\Q)$ preserves whether each of the specified radical coordinates of $\sigma$ is trivial: a conjugation only raises its root-of-unity multiplier to a power prime to the corresponding $s_j$. Chebotarev's theorem supplies infinitely many rational primes $\ell$ with Frobenius in the conjugacy class of $\sigma$, avoiding the prescribed finite set, all listed primes, and $2$. They split completely in $C$, so $\ell\equiv1\pmod m$ and $\ell\equiv1\pmod{s_j}$. At these primes, $q_j$ is an $s_j$th power modulo $\ell$ and every $p\in\mathcal P_j$ is not.

The quotient of $\F_\ell^\times$ by its $s_j$th powers has prime order $s_j$. A nontrivial class stays nontrivial on raising it to any exponent $1\le f\le d<s_j$. Every power of $q_j$, however, is an $s_j$th power. This proves the displayed separation. In particular, overlaps between different lists impose no conflict: their conditions belong to linearly disjoint Kummer extensions with different prime exponents.
\end{proof}

\begin{proposition}\label{prop:same-characteristic}
For all but finitely many primes $w$ of $B$ above rational primes $q$ split in $B$, there exist a prime $v$ of $K$ above the same $q$ and an auxiliary prime $\ell\in\mathcal L$ for which $(v,w)$ is an active pair for $\beta_B$.
\end{proposition}
\begin{proof}
Suppose that there are infinitely many exceptions. Lemma~\ref{lem:finite-fibers} allows us to choose
\[
 w_1,\ldots,w_n,\qquad n>2[K:\Q]+2,
\]
among them, above distinct split rational primes $q_1,\ldots,q_n$, with every $\mathcal V_{w_j}$ finite. Let $\mathcal P_j$ be the set of rational residue characteristics of the primes in $\mathcal V_{w_j}$. The failure of the asserted conclusion means $q_j\notin\mathcal P_j$.

Apply Lemma~\ref{lem:power-separation} with $m=m_0$ and $d=[K:\Q]$. It gives $\ell\in\mathcal L$, away from all the primes involved, such that
\begin{equation}\label{eq:norm-separation}
 Nv\bmod\ell\notin\langle q_j\bmod\ell\rangle
 \qquad(v\in\mathcal V_{w_j},\ 1\le j\le n),
\end{equation}
because $Nv=p^f$ with $p\in\mathcal P_j$ and $1\le f\le[K:\Q]$. If $w_j$ occurred in an active pair at this $\ell$, its source prime would lie in $\mathcal V_{w_j}$. Applying Lemma~\ref{lem:norm-identity} to a lift of source Frobenius and reducing modulo $\ell$ would give
\[
 Nv\bmod\ell\in\langle Nw_j\bmod\ell\rangle
 =\langle q_j\bmod\ell\rangle,
\]
contrary to \eqref{eq:norm-separation}. All $n$ primes $w_j$ therefore belong to $\mathcal M_\ell$, contradicting Proposition~\ref{prop:uniform-missed-primes}.
\end{proof}
